
El plan de implementación define cómo Franco's ERP pasará del diseño conceptual a la operación diaria de Franco's S.A.C. Para ello, organiza el trabajo en fases con hitos verificables, asigna responsables y recursos, define la estrategia de despliegue y reversión, y establece los mecanismos de gestión del cambio y soporte que acompañarán la adopción del sistema. El plan se construye sobre la línea base aprobada en el Entregable 2 (Project Charter, WBS, cronograma, presupuesto y matriz de riesgos), sobre la secuencia de iniciativas de mejora priorizadas en el Entregable 3 y sobre la arquitectura definida en el CE0142, de modo que lo planificado en la gestión del proyecto coincida con lo que se construye en el software y en la infraestructura.

\section{Fases y cronograma}


La implementación corresponde a dos proyectos del Roadmap Tecnológico del Entregable 1 que se ejecutan en paralelo entre el 06/04/2026 y el 30/11/2026 y que atienden el objetivo de TI OET-03 y, a través de él, el objetivo organizacional OEN-03. \textbf{TI-002}, implementación del ERP para la digitalización de procesos, está a cargo del equipo de Desarrollo de Software (SW) y comprende el backend, los dos portales y la base de datos. \textbf{TI-003}, implementación de la infraestructura tecnológica, está a cargo del equipo de Infraestructura Tecnológica (ITI) y comprende la nube, la red de la sede, el respaldo, la continuidad y el monitoreo. Ambos confluyen en el hito H1 del roadmap, previsto para noviembre de 2026: «Base tecnológica operativa: infraestructura organizacional y sistema ERP implementados».

Los dos proyectos dependen entre sí. TI-003 provee el ambiente en el que corre TI-002: la base de datos de producción se aloja en la sede, en dos máquinas virtuales (primaria y en espera), mientras que el backend y los portales se ejecutan en AWS y se comunican con la sede mediante una VPN. Por ello, la red de la sede y el enlace con la nube deben estar operativos antes de las pruebas integrales y de la puesta en producción del sistema.

El trabajo se organiza en siete fases comunes a ambos proyectos y una fase transversal de monitoreo y control. En cada fase, TI-002 y TI-003 desarrollan actividades propias que se coordinan bajo la dirección del Líder del proyecto, y la fase solo se cierra cuando se cumple su hito, como se muestra en la Tabla~\ref{tab:cap4-1}. Al término de la implementación, ambos proyectos continúan con ocho meses de soporte, lo que completa el horizonte de dieciséis meses considerado en la evaluación económica.


{\small
\renewcommand{\arraystretch}{1.3}
\begin{longtable}{|>{\raggedright\arraybackslash}p{1.71cm}|>{\raggedright\arraybackslash}p{1.46cm}|>{\raggedright\arraybackslash}p{3.81cm}|>{\raggedright\arraybackslash}p{3.81cm}|>{\raggedright\arraybackslash}p{3.81cm}|}
\caption{Fases y cronograma de TI-002 y TI-003}\label{tab:cap4-1}\\
\hline
\textbf{Fase} & \textbf{Fechas} & \textbf{TI-002 Sistema ERP (equipo SW)} & \textbf{TI-003 Infraestructura (equipo ITI)} & \textbf{Hito / criterio de cierre} \\
\hline
\endfirsthead
\hline
\textbf{Fase} & \textbf{Fechas} & \textbf{TI-002 Sistema ERP (equipo SW)} & \textbf{TI-003 Infraestructura (equipo ITI)} & \textbf{Hito / criterio de cierre} \\
\hline
\endhead
F-01 Inicio y planificación & 06/04 – 05/05/2026 & Project Charter, plan de gestión y backlog inicial. & Revisión de la infraestructura actual y de las brechas tecnológicas B-08 a B-12. & Project Charter y plan de gestión aprobados. \\
\hline
F-02 Análisis & 06/05 – 25/06/2026 & Requisitos funcionales y no funcionales (SRS), actores del sistema y priorización del backlog. & Requisitos de infraestructura: alojamiento, redundancia, segmentación, respaldo y monitoreo. & SRS validado con Franco's S.A.C. y backlog priorizado. \\
\hline
F-03 Diseño & 25/05 – 25/06/2026 & Arquitectura del monolito modular con Clean Architecture, modelo de datos en PostgreSQL con esquemas por dominio y prototipos de los Portales A y B. & Diseño del esquema sede–nube: datos en la sede, cómputo en AWS, red en VLAN, VPN IPsec, respaldo diario y recuperación con base temporal en la nube. & Arquitectura, modelo de datos y prototipos aprobados. \\
\hline
F-04 Construcción & 06/07 – 23/10/2026 & Ocho sprints que desarrollan los módulos del sistema (detalle en la Tabla~\ref{tab:cap4-2}). & Sede: firewalls pfSense en par, switches en stack con VLAN, dos proveedores de internet y base de datos primaria y en espera sobre VMware. Nube: VPC, VPN, ECS Fargate, CloudFront y S3, WAF, ALB y Secrets Manager. & Módulos funcionales con pruebas unitarias aprobadas. \\
\hline
F-05 Calidad y pruebas & 12/10 – 04/11/2026 & Pruebas unitarias, de extremo a extremo y de aceptación (UAT); integración continua que bloquea cada cambio si fallan la compilación, el lint o las pruebas. & Pruebas de restauración desde el respaldo, de conmutación a la base en espera, de recuperación con la base temporal en la nube y de seguridad del WAF. & $\geq$ 90 \% de casos de prueba aprobados y pipeline CI/CD configurado. \\
\hline
F-06 Despliegue & 02/11 – 27/11/2026 & Puesta en producción de los Portales A y B y del backend. & Ambiente de producción operativo, con monitoreo y alertas activos. & Sistema en producción con disponibilidad $\geq$ 95 \%. \\
\hline
F-07 Cierre & 27/11 – 30/11/2026 & Acta de aceptación, manuales e informe de cierre. & Manual técnico de infraestructura y reporte de disponibilidad. & Acta de aceptación firmada; hito H1 del roadmap alcanzado. \\
\hline
F-08 Monitoreo y control & 05/05 – 27/11/2026 & Seguimiento semanal del avance de los sprints. & Seguimiento de disponibilidad y del gasto en AWS frente al presupuesto mensual. & Informes semanales y reporte quincenal al asesor. \\
\hline
Soporte & 01/12/2026 – 31/07/2027 & Corrección de incidencias y mejoras menores. & Mantenimiento, monitoreo, respaldos y control de costos del ambiente. & Disponibilidad $\geq$ 95 \% durante el soporte. \\
\hline
\end{longtable}
}


La fase de construcción se desarrolla bajo Scrum adaptado, en ocho sprints de dos semanas. Cada sprint entrega un incremento funcional vinculado a las iniciativas de mejora del Entregable 3, lo que permite que el sistema crezca en el mismo orden de prioridad con que se rediseñaron los procesos. La Tabla~\ref{tab:cap4-2} presenta la secuencia.


{\small
\renewcommand{\arraystretch}{1.3}
\begin{longtable}{|>{\raggedright\arraybackslash}p{0.88cm}|>{\raggedright\arraybackslash}p{1.90cm}|>{\raggedright\arraybackslash}p{8.79cm}|>{\raggedright\arraybackslash}p{3.22cm}|}
\caption{Secuencia de sprints de la fase de construcción}\label{tab:cap4-2}\\
\hline
\textbf{Sprint} & \textbf{Fechas} & \textbf{Incremento funcional} & \textbf{Iniciativas (E3)} \\
\hline
\endfirsthead
\hline
\textbf{Sprint} & \textbf{Fechas} & \textbf{Incremento funcional} & \textbf{Iniciativas (E3)} \\
\hline
\endhead
1 & 06/07 – 17/07 & Núcleo transaccional del Portal A: autenticación interna, clientes, productos, ciclo de pedidos y pagos, y control de stock de insumos y productos terminados. & IM-02, IM-05 \\
\hline
2 & 20/07 – 31/07 & Gestión estratégica del negocio: plan y mapa estratégico, KPI e iniciativas con tablero de control. & Gestión estratégica \\
\hline
3 & 03/08 – 14/08 & Gestión estratégica de TI: mapa estratégico, KPI e iniciativas alineados al negocio. & Gestión estratégica TI \\
\hline
4 & 17/08 – 28/08 & Portal B (tienda en línea): registro y verificación del cliente, catálogo, carrito y creación de pedido; configuración base. & IM-01 \\
\hline
5 & 31/08 – 11/09 & Cierre del ciclo de venta y operación de almacén: descuento de inventario y registro de entregas. & IM-01, IM-06 \\
\hline
6 & 14/09 – 25/09 & Trazabilidad entre almacén y producción: solicitudes de insumos y de producción, pedidos en backoffice y notificaciones al cliente. & IM-02, IM-03 \\
\hline
7 & 28/09 – 09/10 & Configuración y backoffice: parámetros del sistema, usuarios internos y flujo pendiente de carrito y pedido del Portal B. & IM-04, IM-05 \\
\hline
8 & 12/10 – 23/10 & Perfil y seguridad de cuenta en ambos portales: datos personales, cambio de contraseña y recuperación de acceso. & Transversal \\
\hline
\end{longtable}
}


La secuencia prioriza primero el núcleo transaccional, porque sobre él se apoyan los demás módulos: no es posible tramitar solicitudes de producción ni medir indicadores sin ventas e inventario registrados en una sola base. Los módulos orientados al cliente se habilitan a partir del sprint 4, cuando el registro interno ya es estable, y los aspectos de seguridad de cuenta se consolidan al final, antes de las pruebas de aceptación.

\section{Responsables y recursos}


La implementación involucra a los integrantes de los equipos de Software e Infraestructura y a dos actores de Franco's S.A.C.: la Gerencia, como patrocinador, y los usuarios clave de las áreas de ventas, inventarios, insumos, producción y despacho, que validan los requisitos y participan en las pruebas de aceptación. La Tabla~\ref{tab:cap4-3} asigna las responsabilidades de cada fase mediante una matriz RACI, en la que R identifica a quien ejecuta la actividad, A a quien la aprueba y rinde cuentas, C a quien se consulta e I a quien se informa; cada actividad tiene un único responsable de aprobación.


{\scriptsize
\renewcommand{\arraystretch}{1.3}
\begin{longtable}{|>{\raggedright\arraybackslash}p{3.00cm}|>{\centering\arraybackslash}p{1.04cm}|>{\centering\arraybackslash}p{1.04cm}|>{\centering\arraybackslash}p{1.04cm}|>{\centering\arraybackslash}p{1.04cm}|>{\centering\arraybackslash}p{1.04cm}|>{\centering\arraybackslash}p{1.04cm}|>{\centering\arraybackslash}p{1.04cm}|>{\centering\arraybackslash}p{1.04cm}|>{\centering\arraybackslash}p{1.04cm}|>{\centering\arraybackslash}p{1.04cm}|}
\caption{Matriz RACI de la implementación}\label{tab:cap4-3}\\
\hline
\textbf{Fase / actividad} & \textbf{Director} & \textbf{Analista} & \textbf{Backend ($\times$2)} & \textbf{DBA y Frontend} & \textbf{Frontend} & \textbf{Diseño y QA} & \textbf{QA y Doc.} & \textbf{Infraestructura ($\times$2)} & \textbf{Gerencia} & \textbf{Usuarios clave} \\
\hline
\endfirsthead
\hline
\textbf{Fase / actividad} & \textbf{Director} & \textbf{Analista} & \textbf{Backend ($\times$2)} & \textbf{DBA y Frontend} & \textbf{Frontend} & \textbf{Diseño y QA} & \textbf{QA y Doc.} & \textbf{Infraestructura ($\times$2)} & \textbf{Gerencia} & \textbf{Usuarios clave} \\
\hline
\endhead
Inicio y planificación & A/R & C & I & I & I & I & R & I & C & — \\
\hline
Análisis & A & R & C & C & I & C & C & C & C & C \\
\hline
Diseño & A & C & R & R & C & R & I & R & I & C \\
\hline
Construcción & A & C & R & R & R & C & I & R & I & I \\
\hline
Calidad y pruebas & A & C & C & C & C & R & R & R & I & C \\
\hline
Despliegue & A & I & C & R & C & C & I & R & C & I \\
\hline
Capacitación de usuarios & A & R & I & I & I & C & R & I & C & C \\
\hline
Cierre & A/R & I & I & I & I & I & R & C & C & — \\
\hline
Soporte post-implementación & A & C & C & C & C & I & C & R & I & C \\
\hline
Control de cambios (transversal) & A & R & C & C & C & I & R & C & C & — \\
\hline
\end{longtable}
}


En cuanto a los recursos, el proyecto dispone del presupuesto aprobado en el Entregable 2 y de los componentes tecnológicos definidos en el CE0142, cada uno con un responsable. Los componentes de la sede aprovechan los equipos disponibles en el laboratorio de la universidad, y los servicios de AWS se controlan con un presupuesto mensual de 80 USD, con avisos al 50, 80 y 100 \% del gasto. La Tabla~\ref{tab:cap4-4} resume los recursos y su control.


{\small
\renewcommand{\arraystretch}{1.3}
\begin{longtable}{|>{\raggedright\arraybackslash}p{2.26cm}|>{\raggedright\arraybackslash}p{4.68cm}|>{\raggedright\arraybackslash}p{1.95cm}|>{\raggedright\arraybackslash}p{4.68cm}|>{\raggedright\arraybackslash}p{1.01cm}|}
\caption{Recursos de la implementación}\label{tab:cap4-4}\\
\hline
\textbf{Tipo} & \textbf{Recurso} & \textbf{Responsable} & \textbf{Uso / control} & \textbf{Costo (S/)} \\
\hline
\endfirsthead
\hline
\textbf{Tipo} & \textbf{Recurso} & \textbf{Responsable} & \textbf{Uso / control} & \textbf{Costo (S/)} \\
\hline
\endhead
Infraestructura (sede) & Servidor con hipervisor VMware para las máquinas virtuales de la base de datos (primaria y en espera) & Equipo de infraestructura & Disponibilidad, réplica y capacidad. & 6,750.00 \\
\hline
Infraestructura (sede y nube) & Licenciamiento VMware y servicios AWS: CloudFront, S3, WAF, ALB, ECS Fargate, ECR, Secrets Manager, VPC y VPN & Equipo de infraestructura & Consumo, disponibilidad y alertas de presupuesto (AWS Budgets). & 20,670.00 \\
\hline
Humano & Especialista cloud, desarrollador y DBA de apoyo & Director del proyecto & Asignación por fase y seguimiento en el tablero. & 2,550.00 \\
\hline
Servicio & Soporte técnico (8 meses $\times$ S/ 2,000) & Equipo de infraestructura & Atención de incidencias y mantenimiento. & 16,000.00 \\
\hline
Red (sede) & Firewalls pfSense en par, switches Cisco C9300L en stack y dos proveedores de internet (Claro y Movistar) & Equipo de infraestructura & Segmentación en VLAN, redundancia de enlace y túneles IPsec hacia AWS. & — \\
\hline
Datos & PostgreSQL 17 con Prisma; volcado diario a S3 y base temporal en RDS ante desastre & DBA & Accesos, integridad, respaldo fuera de la sede y recuperación. & — \\
\hline
Monitoreo & CloudWatch, SNS, Lambda, Step Functions, CloudTrail y Wazuh & Equipo de infraestructura & Detección de la caída de la base, aviso al equipo y registro de actividad. & — \\
\hline
Desarrollo & NestJS, Angular, GitHub y GitHub Actions & Equipo de software & Control de versiones, integración continua y despliegue con credenciales temporales. & — \\
\hline
Gestión & Jira y Figma & Director y Diseñadora & Backlog, avance y prototipos. & — \\
\hline
 & \textbf{Costo total del proyecto} &  &  & \textbf{45,970.00} \\
\hline
\end{longtable}
}


\section{Despliegue y reversión}


\section{Gestión del cambio y soporte}


La gestión del cambio comprende tres frentes: el control de los riesgos de la implementación, la gobernanza que asegura que lo construido corresponda a lo aprobado y la adopción del sistema por parte del personal. A ellos se suma el soporte post-implementación, que sostiene la operación una vez cerrado el proyecto.

\textbf{Riesgos de la implementación.} De la matriz de riesgos del Entregable 2 se seleccionaron los que afectan directamente al despliegue y la adopción del sistema, y se incorporaron dos riesgos que surgen del diseño del CE0142: la dependencia de SICO ERP, que se mantiene como emisor de comprobantes, y el procesamiento de datos en una región de AWS fuera del país. La Tabla~\ref{tab:cap4-5} los presenta con su respuesta.


{\small
\renewcommand{\arraystretch}{1.3}
\begin{longtable}{|>{\raggedright\arraybackslash}p{0.49cm}|>{\raggedright\arraybackslash}p{4.94cm}|>{\raggedright\arraybackslash}p{0.49cm}|>{\raggedright\arraybackslash}p{0.58cm}|>{\raggedright\arraybackslash}p{4.94cm}|>{\raggedright\arraybackslash}p{2.96cm}|}
\caption{Riesgos de la implementación y respuesta}\label{tab:cap4-5}\\
\hline
\textbf{ID} & \textbf{Riesgo} & \textbf{Prob.} & \textbf{Impacto} & \textbf{Respuesta} & \textbf{Responsable} \\
\hline
\endfirsthead
\hline
\textbf{ID} & \textbf{Riesgo} & \textbf{Prob.} & \textbf{Impacto} & \textbf{Respuesta} & \textbf{Responsable} \\
\hline
\endhead
R-01 & Cambios frecuentes de requisitos & 70 \% & Alto & Procedimiento formal de control y aprobación de cambios. & Director del proyecto \\
\hline
R-02 & Retrasos en el desarrollo de los sprints & 50 \% & Alto & Seguimiento semanal, priorización del backlog y reasignación de recursos. & Director / desarrolladores \\
\hline
R-03 & Incremento de ciberamenazas & 50 \% & Alto & Acceso por rol denegado por defecto, WAF delante de los portales y la API, secretos fuera del código y registro de actividad. & Equipo de infraestructura \\
\hline
R-04 & Problemas de integración entre Software e Infraestructura & 45 \% & Medio & Plan de integración y pruebas anticipadas de la conexión entre el backend en AWS y la base de la sede. & Director del proyecto \\
\hline
R-05 & Resistencia al cambio de los usuarios & 45 \% & Medio & Capacitación por rol y acompañamiento durante el despliegue. & QA y Documentadora \\
\hline
R-07 & Incumplimiento de la Ley N.° 29733 & 30 \% & Alto & Protección de datos desde el diseño, cifrado en tránsito y verificación en UAT; consideración del procesamiento en us-east-1. & Analista / desarrollador \\
\hline
R-08 & Fallas en la infraestructura cloud o en la conexión con la sede & 25 \% & Medio & Volcado diario a S3, base temporal en RDS restaurable de forma automática y monitoreo con aviso al equipo. & Equipo de infraestructura \\
\hline
R-09 & Incremento de costos tecnológicos & 25 \% & Bajo & Presupuesto mensual de 80 USD con avisos al 50, 80 y 100 \% y revisión mensual del consumo. & Gerencia / Director \\
\hline
R-10 & Falla de la PC que ejecuta SICO ERP, que no tiene respaldo y sigue emitiendo los comprobantes & — & Alto & Respaldo periódico de la base de SICO ERP y procedimiento alterno de emisión hasta la integración del ERP con SUNAT. & Equipo de infraestructura / Gerencia \\
\hline
\end{longtable}
}


\textbf{Alcance que queda pendiente.} Algunas brechas del CE0141 se atienden solo en parte en esta versión. La Tabla~\ref{tab:cap4-6} registra lo que queda fuera, para que se gestione en las siguientes versiones del sistema.


{\small
\renewcommand{\arraystretch}{1.3}
\begin{longtable}{|>{\raggedright\arraybackslash}p{2.41cm}|>{\raggedright\arraybackslash}p{6.29cm}|>{\raggedright\arraybackslash}p{6.29cm}|}
\caption{Pendientes de las brechas atendidas en parte}\label{tab:cap4-6}\\
\hline
\textbf{Brecha} & \textbf{Qué cubre esta versión} & \textbf{Qué queda pendiente} \\
\hline
\endfirsthead
\hline
\textbf{Brecha} & \textbf{Qué cubre esta versión} & \textbf{Qué queda pendiente} \\
\hline
\endhead
B-03 & Las ventas, el inventario y los lotes envasados pasan del Excel al ERP. & Los documentos de Word no se integran y los archivos históricos no se migran. \\
\hline
B-04 & El ERP envía por correo las verificaciones, credenciales y recuperación de contraseña. & WhatsApp no se integra. \\
\hline
B-09 & Los datos registrados desde la puesta en marcha tienen respaldo fuera de la sede. & Los datos anteriores en papel, Excel y Word no se respaldan con el ERP. \\
\hline
B-14 & Insumos y Producción pasan al ERP. & Compras, Calidad y RR. HH. siguen registrando fuera del sistema. \\
\hline
B-02, B-06, B-07 y B-13 & SICO ERP sigue como emisor oficial de comprobantes. & Integración del ERP con SUNAT y retiro de SICO ERP. \\
\hline
\end{longtable}
}


\textbf{Gobernanza de la implementación.} Ningún cambio en el alcance, el cronograma o el presupuesto se incorpora sin seguir el proceso formal del Entregable 2: solicitud, evaluación de impacto, aprobación del Director y del asesor, implementación y cierre con actualización de la línea base. El Director revisa semanalmente el avance de cada rol e informa al asesor académico de forma quincenal. En el plano técnico, la integración continua impide incorporar cambios al código si fallan la compilación, el lint o las pruebas, y el despliegue hacia AWS se realiza con credenciales temporales, sin claves permanentes en el repositorio.

\textbf{Adopción del sistema.} Para reducir la resistencia al cambio (R-05), la capacitación se organiza por actor del sistema y se realiza en torno a la puesta en producción, de modo que cada usuario aprenda únicamente las funciones que utilizará. La Tabla~\ref{tab:cap4-7} resume el plan.


{\small
\renewcommand{\arraystretch}{1.3}
\begin{longtable}{|>{\raggedright\arraybackslash}p{3.12cm}|>{\raggedright\arraybackslash}p{7.20cm}|>{\raggedright\arraybackslash}p{4.68cm}|}
\caption{Capacitación por actor}\label{tab:cap4-7}\\
\hline
\textbf{Actor} & \textbf{Contenido} & \textbf{Momento} \\
\hline
\endfirsthead
\hline
\textbf{Actor} & \textbf{Contenido} & \textbf{Momento} \\
\hline
\endhead
A1 Operador de Ventas & Clientes, catálogo, registro de pedidos y reembolsos. & Antes de la puesta en producción (F-06) \\
\hline
A2 Operador de Inventarios & Stock de producto terminado, lotes, movimientos y solicitudes de producción. & Antes de la puesta en producción (F-06) \\
\hline
A3 Operador de Insumos & Insumos, sus lotes y movimientos, y atención de solicitudes de insumos. & Antes de la puesta en producción (F-06) \\
\hline
A4 Operador de Producción & Recepción de solicitudes de producción, registro de lotes y envío de solicitudes de insumos. & Antes de la puesta en producción (F-06) \\
\hline
A5 Operador de Despacho & Preparación y entrega de pedidos con sus estados. & Antes de la puesta en producción (F-06) \\
\hline
A6 Analista Estratégico & Objetivos, KPI, iniciativas y tablero del cuadro de mando integral. & Durante la puesta en producción (F-06) \\
\hline
A7 Administrador & Usuarios, roles, empresa emisora y configuración del sistema. & Antes de la puesta en producción (F-06) \\
\hline
A8 Cliente & Registro, verificación de correo, catálogo y cuenta en el Portal B. & Durante la puesta en producción (F-06) \\
\hline
\end{longtable}
}


Esta capacitación inicial cubre el uso del sistema en el arranque. La formación digital más amplia del personal queda a cargo del programa de transformación y capacitación digital (TI-005) previsto en el Roadmap Tecnológico del Entregable 1 para 2027.

\textbf{Soporte post-implementación.} Tras el cierre, el sistema cuenta con ocho meses de soporte técnico a cargo del equipo de infraestructura, con un costo de S/ 2,000 mensuales. El soporte se organiza según las prácticas de ITIL 4: los usuarios reportan las incidencias por un canal único, estas se registran y clasifican por prioridad, y se atienden en el orden que muestra la Tabla~\ref{tab:cap4-8}. Los tiempos son una propuesta del equipo y se validarán con la Gerencia.


{\small
\renewcommand{\arraystretch}{1.3}
\begin{longtable}{|>{\raggedright\arraybackslash}p{1.30cm}|>{\raggedright\arraybackslash}p{7.50cm}|>{\raggedright\arraybackslash}p{6.20cm}|}
\caption{Prioridades de atención en el soporte (propuesta)}\label{tab:cap4-8}\\
\hline
\textbf{Prioridad} & \textbf{Ejemplo} & \textbf{Atención propuesta} \\
\hline
\endfirsthead
\hline
\textbf{Prioridad} & \textbf{Ejemplo} & \textbf{Atención propuesta} \\
\hline
\endhead
Crítica & El sistema no permite registrar ventas o inventario. & Inicio inmediato; solución en el mismo día. \\
\hline
Alta & Falla en un módulo con alternativa temporal. & Dentro de las 24 horas. \\
\hline
Media & Error que no detiene la operación. & Dentro de las 72 horas. \\
\hline
Baja & Consulta de uso o mejora menor. & Se programa en el siguiente mantenimiento. \\
\hline
\end{longtable}
}


Durante este periodo se mantienen activos los mecanismos de operación definidos en el CE0142: una sonda prueba cada minuto la conexión con la base de datos y, si detecta una caída, avisa al equipo y activa la recuperación con la base temporal en la nube; el volcado diario de la base se conserva en S3 fuera de la sede; y el registro de actividad de la cuenta de AWS se revisa desde Wazuh. La disponibilidad se registra cada mes para verificar la meta del 95 \%, y las incidencias recurrentes y las mejoras solicitadas se canalizan mediante el proceso de control de cambios, lo que conecta el soporte con el seguimiento de indicadores que se desarrolla en el capítulo 5.
